% Include from a report located in reports/, after running doit part4.
\input{../_output/part4/results_macros.tex}
\section{Maturity effects and spot/roll decomposition}
We replicate the supplied Tables IA.XI and IA.XV (course Tables 4 and 5)
using WRDS Datastream and a 1986--2014 historical evaluation window.
Contract-specific missing histories and a shorter sample prevent exact
comparison with the authors' 1960--2014 CRB/FII results.

\begin{table}[htbp]
\centering
\caption{Basis-momentum across the curve. The table compares equal-weight
High4-minus-Low4 returns at different nearby maturities. Conditional rows use
same-side portfolio agreement of at most three commodities; unsigned-overlap
sensitivity results are exported separately. Annual means are in percent;
t statistics use monthly iid standard errors. Per-cell sample counts differ
when selected contracts cannot be repriced.}
\resizebox{\textwidth}{!}{\input{../_output/part4/table4.tex}}
\end{table}

\begin{table}[htbp]
\centering
\caption{Spot and roll decomposition. In the substituted WRDS sample,
the basis-momentum return differential is associated primarily with the roll
component; the spot differential is negative. Components are calculated
per commodity before averaging, with identical evaluation months across
groups. The spot component is a residual rather than observed physical spot
returns. The multiplicative interaction is retained in the audit outputs.}
\resizebox{\textwidth}{!}{\input{../_output/part4/table5.tex}}
\end{table}

The annualized basis-momentum High4-minus-Low4 futures return is
\PartFourFutures\% ($t=\PartFourFuturesT$), compared with a roll component of
\PartFourRoll\% and a spot component of \PartFourSpot\%.
These are historical gross returns, not an updated or transaction-cost-adjusted
investment result. Gasoline coverage ends in 2006, and anomalous contract
changes in rough rice and silver are excluded from the decomposition.
